\documentclass[leqno, a4paper, 12pt]{jsarticle}
\usepackage{amssymb}
\usepackage{amsthm}
\usepackage{amsmath}
\usepackage{comment}
\usepackage{url}

\usepackage{enumitem}
%\usepackage{showkeys}


%定理環境 thmはあまり使わずpropをなるべく使う。
%主張は基本的に証明の中で使い、そのため番号を付けない。
%注意環境を付けてみた。
\newtheorem{thm}{定理}
\newtheorem{prop}[thm]{命題}
\newtheorem{cor}[thm]{系}
\newtheorem{lem}[thm]{補題}
\newtheorem{df}[thm]{定義}
\newtheorem{exam}[thm]{例}
\newtheorem{fact}[thm]{事実}
\newtheorem*{claim}{主張}
\newtheorem*{cau}{注意}
\newtheorem*{rmk}{注意}
\renewcommand{\proofname}{証明}
%矢印たち。
%\toは\rightarrowと同じ。\getsは\leftarrowと同じ。同値は\iff
\newcommand{\To}{\Longrightarrow}
\newcommand{\Gets}{\LongLeftarrow}
%黒板太字
\newcommand{\nn}{\mathbb{N}}
\newcommand{\zz}{\mathbb{Z}}
\newcommand{\qq}{\mathbb{Q}}
\newcommand{\rr}{\mathbb{R}}
\newcommand{\cc}{\mathbb{C}}
%無限大
\newcommand{\mgn}{\infty}
%差集合は\setminusで統一する。等号の否定は\neq
%真部分集合は\subsetneqq　偏微分は\partial
%ディスプレイスタイル。\dfracでディスプレイ分数
\newcommand{\dpst}{\displaystyle}

\title{論文メモ
\large{\\--- なにこれ？ ---}
}
\author{ナンブ　キトラ}
\date{\today}
\begin{document}
\maketitle

本棚を整理したいので，
溜まっている論文のメモを書いて未来への手紙とすることにする．

なぜこれらの論文が同じ束に入っていたのかは謎である．
もしかするとFund. Math.のサイトを見て66年くらいの論文で面白そうなのを適当に印刷したのかもしれない．

論文
\cite{MR0202117}
では逆極限空間の写像が近似できるとかホモトピックとか
そういう連続体論的な話をしているようだ．



論文
\cite{MR0200910}
は
ちょっとよくわかんなかった．
Moore decomposition 
の話が載っていて，
これは特殊なUpper semi-continuous decomposition (USCD)である
(ケリーにUSCD載ってるよ．今ふうにいうと確か閉同値関係による分割だったと思う)．
こういうのはなんか見たことがあるので
そこら辺からこの論文に辿り着いたのかもしれない．


論文
\cite{MR0180956}
は有名な長田の$n$次元距離化可能空間の特徴づけというやつで，距離化可能空間$X$が$\dim X\le n$であるための必要十分条件は$X$の位相と両立する距離関数$\rho$があって，
任意の$n+3$点$x, y_{1}, \dots, y_{n}, y_{n+1}$
に対してある異なる番号$i, j$が存在して
$\rho(y_{i}, y_{j})\le \rho(x, y_{j})$
を満たすことであると証明している．
この定理の前身として長田自身の似たような定理や
de Groot の可分な場合の上の定理などがあり，
それらを踏まえて一般的に証明したようだ．
論文
\cite{MR2666207}
と
\cite{MR1906836}
も参照のこと．
さらに論文
\cite{MR0982791}
\cite{MR0086284}
\cite{MR0180957}
\cite{MR0164320}
なども参考に．





%READ!
%myplain is the same to amsplain style. 
\nocite{*}
\bibliographystyle{myplaindoidoi}
\bibliography{hatena.bib}



\end{document}